\section{Conclusion}
\label{sec:cn-conclusion}

This project presented \textbf{L2M}, an LLM-driven lyrics-to-melody
generation system that achieves syllabic alignment, emotional coherence, and
operational reliability without any task-specific training. The core
technical insight is a two-stage prompting strategy --- emotion analysis
followed by conditioned melody generation --- grounded in an explicit
emotion-to-musical-attribute mapping, with a deterministic fallback layer
that guarantees pipeline completion for every input.

Quantitative evaluation on a 20-sample benchmark spanning six emotion
categories demonstrated 100\% syllable--note alignment, 95\% emotion--key
consistency, 100\% MIDI validity, and an average end-to-end latency of 3.2
seconds. Qualitative review confirmed that generated melodies are diatonic,
singable, and emotionally appropriate to their lyrics. The system is
released as an open-source Python package with a command-line interface and
programmatic API, providing an accessible, reproducible baseline for
LLM-based music generation research.

Returning to the research questions of \Autoref{chap:intro}: a
general-purpose LLM \emph{can} generate structurally valid, emotionally
consistent melodies zero-shot when its output is constrained by structured
prompts and typed validation (RQ1); two-stage decomposition improves both
interpretability and fallback precision (RQ2); a deterministic
music-theory-grounded fallback converts a 15\% LLM failure rate into a 100\%
system completion rate at an acceptable quality cost (RQ3); and hard
constraints, output schemas, and few-shot examples are the prompt
engineering techniques that make the difference between unusable and
reliable output (RQ4). More broadly, the results suggest that zero-shot LLM
prompting, carefully engineered with structured constraints, is a viable and
practical alternative to supervised training for constrained creative
generation tasks --- particularly where training data is scarce and
reliability is a requirement.

\section{Future Work}
\label{sec:cn-future}

Five directions extend this work naturally.

\subsection{Harmonic Generation}

The current system produces monophonic melodies. Future versions could
generate chord progressions conditioned on the melodic output, support
multi-voice arrangements (harmony, bass line, countermelody), and perform
automatic instrumentation and orchestration.

\subsection{Improved LLM Integration}

Larger proprietary models and open-source alternatives should be evaluated
for quality--cost--reliability trade-offs, including locally hosted models
that would remove the API dependency identified in
\Autoref{sec:ev-limitations}. A compact LLM fine-tuned on paired
lyrics--melody data --- retaining the structural prompt framework ---
could improve creative quality, and reinforcement learning from human
feedback using musicologist-rated outputs offers a principled path to
optimizing musicality directly.

\subsection{Formal Evaluation}

A large-scale user study (musicians and general listeners, $N \geq 100$)
using Mean Opinion Scores and A/B comparison against systems such as ReLyMe
\cite{zhang22relyme} and SongComposer \cite{ding24songcomposer} would
substantiate the qualitative findings. Complementary musicological metrics
--- pitch-class distribution entropy, interval histogram similarity to
professional compositions, and phrase boundary detection accuracy --- would
enable objective comparison at scale.

\subsection{Non-Western Music Support}

The emotion--key mapping could be extended to non-Western systems, including
maqam (Arabic), raga (Indian), and pentatonic (East Asian) frameworks,
together with multi-language syllabification for non-English lyrics ---
addressing the cultural scope limitation directly.

\subsection{Interactive Generation}

A web-based interface with real-time playback and in-browser melody editing
would open the system to non-programmer musicians. An iterative refinement
loop --- in which the user adjusts the detected emotion label or selected
key and the system regenerates --- would exploit the interpretability
advantage of the two-stage design.

\vspace{1em}
\noindent
In closing, this work demonstrates that the reasoning capabilities of
modern language models, when disciplined by structured prompting, typed
validation, and music-theoretic grounding, are sufficient to build a
practical, reliable bridge from words to music.
